\section{Обзор литературы}

    \subsection{Введение в кластеризацию}

        Кластеризация представляет собой один из фундаментальных методов машинного обучения, используемый для группировки объектов на основе их внутренних признаков и сходства без привлечения учителя, то есть без использования размеченных данных. В отличие от классификации, где объекты распределяются по заранее определенным классам, кластеризация позволяет выявлять скрытые структуры и закономерности в данных, что делает ее особенно полезной в ситуациях, когда заранее неизвестно, как должны выглядеть классы или группы.

        Кластеризация применяется в широком спектре научных и прикладных задач, охватывая такие области, как анализ данных \cite{SETYANINGSIH2012286}, компьютерное зрение \cite{Demidovskij2023ids}, биоинформатика \cite{Kiselev2019scRNAseqClustering}, обработка естественного языка \cite{Grootendorst2022BERTopic}, а также информационные технологии, социальные науки, медицина и финансовая аналитика. Ее основная цель -- разделить множество объектов на группы (кластеры) таким образом, чтобы объекты в пределах одного кластера обладали максимальной степенью сходства, в то время как объекты различных кластеров были существенно различны.

        Кластеризация может быть использована для решения различных задач, таких как:
        \begin{itemize}
            \item \textbf{Сегментация изображений}: разделение изображения на области схожих пикселей, что позволяет выделить объекты на изображении и улучшить качество последующего анализа.
            \item \textbf{Классификация текстов}: группировка текстовых документов по темам или стилям, что позволяет упростить поиск и анализ информации.
            \item \textbf{Кластеризация пользователей}: сегментация пользователей на группы схожих интересов и предпочтений, что позволяет персонализировать рекомендации и маркетинговые стратегии.
            \item \textbf{Обнаружение аномалий}: выявление аномальных точек, не принадлежащих ни одному кластеру или попадающих в редкие группы, что может свидетельствовать о мошенничестве, сбоях в оборудовании или ошибках в данных.
            \item \textbf{Разведочный анализ данных (EDA, Exploratory Data Analysis)}: выявление скрытых закономерностей и структур в данных, что особенно полезно при работе с новыми или плохо изученными наборами данных.
        \end{itemize}

    \subsection{Методы кластеризации}
        Существуют различные подходы к реализации кластеризации, отличающиеся не только алгоритмической природой, но и способностью справляться с различными типами данных и структурами кластеров. В зависимости от метода формирования кластеров можно выделить следующие основные категории:

        \begin{itemize}
            \item \textbf{Центроидные методы} -- основаны на минимизации расстояний между объектами и центроидами (центральными точками) кластеров (\kmeans{}, \kmedoids{}).
            \item \textbf{Иерархические методы} -- формируют древовидную структуру кластеров путем последовательного объединения или разбиения групп (агломерационная и дивизионная кластеризация).
            \item \textbf{Плотностные методы} -- строят кластеры как области повышенной плотности объектов (например, DBSCAN, OPTICS).
            \item \textbf{Модельные методы} -- предполагают, что данные генерируются из смеси вероятностных распределений (например, EM-алгоритм для гауссовых смесей).
            \item \textbf{Спектральные и графовые методы} -- используют спектральные свойства графов и матриц смежности для обнаружения кластеров в данных со сложной топологией.
        \end{itemize}

        Каждый из подходов обладает своими достоинствами и ограничениями, что требует осознанного выбора метода в зависимости от конкретных характеристик исследуемых данных: размерности, плотности, формы кластеров и наличия шумов.

        \subsubsection{Кластеризация на основе центроидов}

            Кластеризация на основе центроидов (centroid-based clustering) представляет собой один из наиболее интуитивно понятных и широко используемых подходов к группировке данных. Суть данного метода заключается в том, что каждый кластер описывается так называемым центроидом -- обобщенным представлением <<центра тяжести>> данных в пределах соответствующего кластера. Объекты присваиваются к тому или иному кластеру на основе минимизации расстояния до соответствующего центроида.

            Наиболее известным и базовым представителем  кластеризации на основе центроидов является алгоритм \kmeans{} \cite{Lloyd1982kmeans}. Данный алгоритм итеративно обновляет позиции центроидов и перераспределяет объекты между кластерами с целью минимизации инерции -- суммарного квадратичного расстояния между объектами и центрами кластеров. Простота реализации и высокая вычислительная эффективность обеспечили \kmeans{} широкое распространение, особенно при работе с большими объемами данных.

            Однако классический \kmeans{} обладает рядом существенных ограничений. Во-первых, он чувствителен к начальной инициализации центроидов, что может приводить к сходимости к локальным минимумам. Во-вторых, алгоритм предполагает заданное заранее число кластеров \(k\), что не всегда возможно определить. Кроме того, алгоритм плохо справляется с кластерами, имеющими сложную форму, различную плотность или отличающиеся размерами.

            В ответ на указанные ограничения было предложено несколько модификаций. Один из таких подходов -- \kmeans{}++ \cite{Arthur2010kmeans++}, который улучшает этап начальной инициализации центроидов путем применения стохастической стратегии выбора начальных точек. Это повышает стабильность и качество итогового разбиения, особенно при наличии шумов и неоднородности в данных.

            Альтернативой, сохраняющей идею центроидного представления, является алгоритм \kmedoids{}, также известный как PAM (Partitioning Around Medoids) \cite{kaufman2008kmedioids}. В отличие от \kmeans{}, где центроиды могут быть абстрактными точками, \kmedoids{} использует в качестве центров реальные объекты из выборки -- медианы. Это повышает устойчивость к выбросам и позволяет использовать произвольные метрики расстояния, включая манхэттенскую, косинусную и другие, однако вычислительная сложность алгоритма существенно выше, особенно при больших объемах данных.

            Среди более современных подходов можно выделить CLARANS (Clustering Large Applications based upon RANdomized Search) \cite{Ng1994clarans}, который комбинирует элементы \kmedoids{} с методами стохастического поиска и улучшает масштабируемость алгоритма при работе с крупными выборками. В свою очередь, Mini-Batch \kmeans{} \cite{Sculley2010minibatch} адаптирован для обработки потоков данных и масштабных задач за счет обновления центроидов на основе случайных подвыборок (батчей), что уменьшает требования к памяти и ускоряет обучение.

            Несмотря на ограниченность в обнаружении кластеров сложной формы, методы, основанные на центроидах, остаются актуальными благодаря своей простоте, интерпретируемости и высокой вычислительной эффективности. Они особенно полезны в качестве базовых моделей или в составе более сложных гибридных решений, где предварительная кластеризация используется как этап инициализации или предобработки данных.

        \subsubsection{Кластеризация на основе плотности}

            Методы кластеризации на основе плотности (density-based clustering) представляют собой эффективный класс алгоритмов, разработанных для выявления кластеров в данных с неоднородной плотностью распределения. Основная идея заключается в том, что кластеры формируются как связные области пространства с высокой плотностью объектов, отделенные друг от друга зонами разреженности. В отличие от методов, основанных на центроидах или иерархической структуре, плотностные алгоритмы способны выделять кластеры произвольной формы, устойчивы к шуму и выбросам, а также не требуют предварительного задания числа кластеров.

            Наиболее известным и фундаментальным методом в этой категории является DBSCAN (Density-Based Spatial Clustering of Applications with Noise) \cite{Ester1996DBSCAN}. Алгоритм определяет кластеры как максимальные области, содержащие взаимосвязанные точки с плотностью выше заданного порога. Основные параметрами алгоритма радиус окрестности и минимальное число точек в окрестности позволяют определить является ли точка ядровой, граничной или шумовой. DBSCAN эффективно справляется с задачами обнаружения кластеров различных форм, но его производительность существенно зависит от корректного выбора параметров, что особенно проблематично при переменной плотности данных.

            Для преодоления ограничений DBSCAN в условиях разной плотности кластеров был предложен алгоритм OPTICS (Ordering Points To Identify the Clustering Structure) \cite{Ankerst1999OPTICS}. В отличие от DBSCAN, OPTICS не производит окончательное разбиение, а строит упорядоченное представление плотности данных, называемое reachability plot. Это позволяет визуально определить плотностные структуры и формировать кластеры постфактум, адаптируясь к локальным особенностям плотности. OPTICS расширяет применимость методов плотностной кластеризации, особенно в случаях с вложенными и слабо выраженными структурами.

            Более современным и мощным развитием идей DBSCAN является HDBSCAN (Hierarchical DBSCAN) \cite{Maarten2022HDBSCAN}. Данный алгоритм строит иерархическую дендрограмму на основе плотностной связи между точками, а затем применяет процедуру извлечения стабильных кластеров. Он автоматически определяет оптимальное количество кластеров и значительно лучше справляется с кластерами переменной плотности, при этом наследуя устойчивость DBSCAN к шуму. HDBSCAN также предлагает оценку надежности принадлежности точки к кластеру, что повышает интерпретируемость результатов.

            Другим интересным направлением является DENCLUE \cite{DENCLUE}, использующий ядерные оценки плотности (kernel density estimation) для моделирования распределения данных. Это позволяет формировать кластеры на основе градиентных потоков в плотностном пространстве и обеспечивает гибкую настройку формы и гладкости кластеров. Однако метод требует вычислительно дорогих операций и чувствителен к выбору параметра сглаживания.

            Плотностные алгоритмы находят широкое применение в задачах обнаружения аномалий, анализе геопространственных данных, выявлении сообществ в социальных сетях и анализе временных рядов, где традиционные методы центроидного типа или иерархические подходы оказываются неэффективными.

            В целом, подходы, основанные на плотности, являются мощными средствами анализа сложных, неструктурированных и шумных данных. Их гибкость и устойчивость делают их незаменимыми в случаях, когда структура кластеров заранее неизвестна и может отличаться высокой степенью нелинейности.

        \subsubsection{Иерархическая кластеризация}
            Иерархическая кластеризация представляет собой мощный класс алгоритмов, ориентированных на поэтапное построение структуры вложенных кластеров без необходимости заранее задавать их количество. В отличие от центроидных и плотностных методов, иерархический подход формирует древовидную структуру, называемую дендрограммой, где каждый уровень отражает определенный этап объединения или разбиения кластеров. Это делает метод особенно полезным для визуального анализа структуры данных и выбора оптимального количества кластеров после завершения обучения.

            Существует два основных направления иерархической кластеризации:

            \begin{itemize}
                \item Агломеративная стратегия, при которой каждый объект изначально рассматривается как отдельный кластер, и далее происходит последовательное объединение ближайших кластеров вплоть до образования единой группы. Данный подход является наиболее распространенным и практически применимым, особенно в случае умеренных размеров выборки. Одним из самых популярных алгоритмов этой категории является метод Уорда \cite{Murtagh2011wardalgo}, минимизирующий увеличение внутрикластерной дисперсии при каждом объединении.
                \item Дивизионная стратегия, напротив, начинается с единственного кластера, включающего все объекты, и рекурсивно разделяет его на подмножества. Несмотря на высокую концептуальную привлекательность, методы этого типа требуют значительно больших вычислительных затрат и применяются реже \cite{Kaufman1990divisive}.
            \end{itemize}

            Качество работы иерархической кластеризации в значительной степени зависит от выбора метрики расстояния между кластерами. Существует несколько стратегий, определяющих, как именно измеряется расстояние между группами данных:

        \paragraph{Single Linkage (Метод ближнего соседа)}
            Метод ближнего соседа \cite{Sibson1973SLINK} определяет расстояние между кластерами как минимальное расстояние между парой точек из разных кластеров. Он хорошо справляется с кластеризацией вытянутых или нестандартных форм, однако чувствителен к шуму и может приводить к так называемому «эффекту цепочки» -- последовательному объединению разреженных объектов.

        \paragraph{Complete Linkage (Метод дальнего соседа)}
            В методе дальнего соседа \cite{Defays1977CLINK} расстояние между кластерами определяется как максимальное расстояние между точками двух кластеров. Этот подход обеспечивает более компактные и равномерные кластеры, но при этом склонен переоценивать расстояния и может неправомерно разделять вытянутые структуры.

        \paragraph{Average Linkage (Метод средней связи)}
            Средняя связь \cite{Finkelman2021averagelinkage} использует среднее значение расстояний между всеми парами точек из двух кластеров. Данный метод обеспечивает баланс между крайними подходами и демонстрирует устойчивость к шуму, но иногда недостаточно точно отражает геометрию данных.

        Кроме того, иерархические методы удобны тем, что позволяют динамически выбирать оптимальное число кластеров, обрезая дендрограмму на нужной высоте без необходимости многократного запуска алгоритма. Это особенно ценно в условиях, когда структура данных неизвестна заранее.

        Несмотря на высокую интерпретируемость и визуальную наглядность, недостатками иерархических методов являются высокая вычислительная сложность ( \(\mathcal{O}(n^2)\) или выше для большинства реализаций) и невозможность коррекции ошибок: однажды принятые решения о слиянии или разбиении не подлежат пересмотру.

        Тем не менее, иерархическая кластеризация остается важным инструментом в арсенале анализа данных и активно применяется в биоинформатике, текстовом анализе, социологии, психологии и других дисциплинах, где требуется исследование сложных иерархических структур в выборке.

    \subsubsection{Кластеризация на основе обучения Хебба}
        Кластеризация на основе обучения Хебба (Hebbian Learning) представляет биологически вдохновленный метод обучения, основывающийся на фундаментальном принципе нейрофизиологии. Согласно так называемому правилу Хебба \cite{Hebb1949}, «нейроны, которые активируются совместно, укрепляют связь между собой» (neurons that fire together, wire together). Эта идея легла в основу адаптивного обучения в искусственных нейронных сетях, а также стала мощным инструментом для кластеризации данных на основе корреляционной близости объектов.

        С точки зрения формализации, классическое обучение по правилу Хебба предполагает обновление весов связей по правилу:
        \begin{equation}
            \Delta w = \eta x y,
        \end{equation}
        где \(x\) и \(y\) -- входной и выходной сигналы, \(w\) -- вес связи между входом и выходом, а \(\eta\) -- коэффициент обучения. При многократном предъявлении данных такие схемы обучения приводят к усилению связей между часто совместно активируемыми признаками, что способствует формированию естественных кластеров.

        Преимуществом данного метода кластеризации является биологическая правдоподобность, адаптивность к данным, способность выявлять латентные структуры без строгих априорных предположений о количестве кластеров. Однако они также обладают определенными ограничениями. Наиболее существенными из них являются чувствительность к параметру обучения и зависимость от порядка предъявления входов.

        Тем не менее, обучение по правилу Хебба продолжает рассматриваться как перспективная альтернатива классическим методам кластеризации, особенно в условиях отсутствия четких границ между группами данных.


    \subsection{Проблемы пространства большой размерности}
        Анализ и кластеризация данных в пространствах высокой размерности остаются актуальными задачами современной науки о данных. По мере увеличения числа признаков (измерений), используемых для описания объектов, традиционные методы анализа, включая алгоритмы кластеризации, начинают сталкиваться с фундаментальными затруднениями. Это явление широко известно как <<проклятие размерности>> (curse of dimensionality) \cite{Bellman1957cursedim}.

        \paragraph{Вычислительная сложность}
            Одним из наиболее очевидных последствий высокой размерности является экспоненциальный рост вычислительной нагрузки. Многие алгоритмы кластеризации, в частности \kmeans{}, требуют вычисления расстояний между каждым объектом и каждым центроидом на каждой итерации. При этом сложность возрастает линейно по числу измерений и квадратично по числу объектов. В пространствах с тысячами признаков это приводит к существенному увеличению времени работы алгоритмов и ресурсов памяти, ограничивая масштабируемость \cite{Verleysen2005curse}.

        \paragraph{Разреженность данных}
            По мере увеличения размерности пространства наблюдается феномен разреженности -- объекты становятся статистически эквивалентно удаленными друг от друга, и плотность данных резко падает. Как показано в \cite{Beyer1999NearestNeighbor}, различия в расстояниях между ближайшими и дальними соседями стремятся к нулю по мере роста размерности, делая евклидову метрику все менее информативной \cite{aggarwal2001surprising}. Это приводит к затруднениям в определении границ между кластерами, особенно при использовании методов на основе расстояния и плотности (например, \kmeans{}, DBSCAN).

    \subsubsection{Использование расстояния Минковского}
        Одной из возможных стратегий повышения качества кластеризации в пространствах большой размерности является адаптация метрики расстояния. В этом контексте особое внимание заслуживает расстояние Минковского, которое представляет собой обобщение различных метрических функций:
        \begin{equation}
            \label{eq:minkowski}
            d(x, y) = \left( \sum_{i=1}^{n} |x_i - y_i|^p \right)^\frac{1}{p}.
        \end{equation}
        Значения параметра \(p\) позволяют интерполировать между манхэттенской (\(p=1\)) и евклидовой (\(p=2\)) метриками, а при \(p \to \infty\) метрика приближается к максимуму по координатам. Исследования показали, что выбор оптимального значения параметра \(p\) может существенно повлиять на эффективность кластеризации в условиях высокой размерности \cite{aggarwal2001surprising}. Особенно перспективным направлением является использование взвешенного расстояния Минковского, где каждый признак имеет индивидуальный вес:
        \begin{equation}
            \label{eq:weighted_minkowski}
            d(x, y) = \left( \sum_{i=1}^{n} w_i |x_i - y_i|^p \right)^\frac{1}{p},
        \end{equation}
        где \(w_i\) -- вес признака \(i\). Эти веса могут быть получены, например, с помощью анализа важности признаков, извлеченной из моделей машинного обучения, либо с использованием разложения главных компонент (PCA). Такой подход позволяет учитывать разную информативность признаков, что особенно важно в условиях высокой размерности и коррелированных признаков. Практические результаты подтверждают, что взвешенные метрики существенно улучшают устойчивость алгоритмов, таких как \kmeans{}, к шуму и нерелевантным признакам \cite{chouikhi2015featureweighting}.

    \subsubsection{Кластеризация с частичным привлечением учителя}
        Другая важная стратегия повышения эффективности кластеризации -- это частичное привлечение учителя (semi-supervised clustering). Данная стратегия сочетает преимущества обучаемых без учителя методов с ограниченным числом меток классов, что позволяет направлять процесс кластеризации к более интерпретируемым и структурированным результатам. Такой подход используется, например, при наличии  информации о небольшом числе объектов, чья принадлежность к кластерам заранее известна, что позволяет добавить контроль в процесс кластеризации, приводя к созданию различимых кластеров с более четкой структурой, а также добавляя возможность использовать стандартные методы оценки качества результатов алгоритма, такие как точность (Accuracy), прецизионность (Precision) полнота (Recall). Среди алгоритмов кластеризации с частичным привлечением учителя можно выделить основные подходы использования ограниченного знания о классах исходных данных \cite{hse:5_course}.

        \paragraph{Привлечение учителя для генерации начальных кластеров}
            Подход кластеризации с частичным привлечением учителя для генерации начального разбиения на кластеры (seeding) \cite{basu2002seeding}, \cite{antoine2014seeding} -- это направление, в котором небольшой набор маркированных точек данных используют для формирования начальных кластеров, таким образом предоставляя алгоритму хорошую стартовую точку для оптимизационной задачи. В дальнейшем алгоритм кластеризации уже без учителя распространит начальные метки на оставшиеся немаркированные данные на основе мер схожести, обеспечивая сильное влияние начального приближения на весь процесс обучения.

        \paragraph{Привлечение учителя для определения принадлежности к кластеру}
            Одним из способов внести знания о кластерах в процесс обучения также является привлечение учителя для определения принадлежности к кластеру. Метод Semi-supervised clustering under a <<compact-cluster>> assumption \cite{jiang2023compactcluster}, который основан на допущении идеи компактности кластеров и минимизирует внутрикластерные расстояния, одновременно максимизируя межкластерные расстояния, обеспечивает четкость и компактность кластеров.

    \subsubsection{Self-Supervised Clustering}
        Самоконтролируемая кластеризация (Self-Supervised Clustering, SSL) представляет собой гибридную парадигму, сочетающую принципы самоконтролируемого обучения (Self-Supervised Clustering) с задачей кластеризации. В ее основе лежит идея извлечения устойчивых и структурно значимых признаков из неразмеченных данных, которые в дальнейшем могут использоваться для формирования кластеров с высокой интерпретируемостью и разделимостью. SSL зарекомендовал себя как эффективный инструмент предварительного анализа, способный существенно улучшать производительность задач, включая кластеризацию \cite{Gui2023SSL}.

        SSL представляет собой класс методов, в которых обучающая информация извлекается из самих данных без использования явных меток. В результате обучения создаются эмбеддинги (embeddings) -- векторы признаков, которые можно использовать в задачах сегментации, визуализации и кластеризации.

        Таким образом, SSL становится этапом предварительного анализа данных. Одним из таких методов может быть этап снижения размерности, позволяющий компактно представить данные в пространстве с меньшей размерностью, сохранив при этом ключевую топологическую структуру.

    \subsubsection{Методы снижения размерности}

        \paragraph{Метод главных компонент (Principal Component Analysis)}
            Метод главных компонент (PCA) -- это один из наиболее известных и широко используемых алгоритмов снижения размерности \cite{Jolliffe1986pca}. PCA позволяет перевести данные из пространства высокой размерности в подпространство меньшей размерности, сохранив при этом максимум возможной дисперсии данных. Важной особенностью метода является его линейная природа, что делает его особенно эффективным для задач, в которых основная информация содержится в линейных комбинациях признаков. Простота PCA, его высокая вычислительная эффективность и математическая интерпретируемость сделали метод фундаментальным инструментом в таких областях, как:
            \begin{itemize}
                \item обработка изображений (например, сжатие изображений, распознавание лиц \cite{Turk1991faces}),
                \item биоинформатика (анализ экспрессии генов \cite{Ringner2008pca}),
                \item финансовое моделирование (анализ корреляционных структур рынков),
                \item кластеризация и визуализация многомерных данных \cite{Jolliffe2016PCAbook}.
            \end{itemize}
            Однако, несмотря на достоинства, линейность PCA является его ограничением: метод не способен корректно обрабатывать данные, в которых основная структура отражается в нелинейных зависимостях между признаками.

        \paragraph{Метод многомерного шкалирования (Multidimensional Scaling)}
            Метод многомерного шкалирования (MDS) -- это классический подход к снижению размерности, направленный на сохранение попарных расстояний или близостей между объектами при отображении их в пространство меньшей размерности. В отличие от метода главных компонент (PCA), MDS может использовать любой тип метрик, включая неевклидовы, и поэтому считается более гибким для анализа данных с нелинейной структурой.

            MDS был впервые разработан как способ визуализации психологической близости между объектами на основе субъективных оценок \cite{Kruskal1964mds}. Позднее метод получил широкое распространение в информационных науках, биоинформатике, социальной аналитике, анализе текстов и других дисциплинах, где структура данных может быть выражена через матрицу попарных расстояний.

        \paragraph{Метод Uniform Manifold Approximation and Projection (UMAP)}
            Метод UMAP \cite{McInnes2018UMAP} представляет собой современный и эффективный нелинейный алгоритм снижения размерности, ориентированный на сохранение локальной и глобальной структуры данных. В отличие от классических линейных методов (например, PCA), UMAP обладает высокой масштабируемостью, вычислительной эффективностью и устойчивостью к шуму, что делает его востребованным в задачах визуализации, кластеризации и предварительной обработки данных.

        \paragraph{Метод Isometric Mapping (ISOMAP)}
            Метод ISOMAP -- это нелинейный алгоритм снижения размерности \cite{Tenenbaum2000isomap}, расширяющий классический метод многомерного шкалирования (MDS) и сохраняющий геодезические расстояния между точками, а не только их евклидову дистанцию. ISOMAP позволяет эффективно обрабатывать данные с нелинейной структурой, сохраняя при этом важные топологические свойства. Однако метод требует значительных вычислительных ресурсов и может быть чувствителен к шуму в данных.

        \paragraph{Метод Locally Linear Embeddings (LLE)}
            Метод LLE -- это эффективный алгоритм нелинейного снижения размерности \cite{Roweis2000LLE}, позволяющий получить компактное представление высокоразмерных данных, сохранив при этом локальную геометрию исходного пространства. В отличие от линейных методов, таких как PCA, LLE способен выявлять нелинейные структуры, что делает его особенно полезным для обработки данных, лежащих на искривленных многообразиях.

        \paragraph{Метод Spectral Embeddings (SSE)}
            Метод SSE, также известный как спектральное разложение графа или Laplacian Eigenmaps, представляет собой мощный инструмент нелинейного снижения размерности \cite{Belkin2001laplacian}. SSE опирается на представление данных в виде графа, в котором вершины соответствуют объектам, а ребра -- их близости. Метод позволяет сохранить локальную геометрию данных, что делает его особенно эффективным для задач, связанных с визуализацией, сегментацией и кластеризацией структурно сложных данных.

    \subsection{Выводы}
        В ходе проведенного обзора были рассмотрены основные алгоритмы кластеризации, применяемые в задачах анализа многомерных данных. Одной из ключевых проблем, выявленных при этом, стало существенное снижение качества кластеризации в условиях высокой размерности данных. Такая ситуация, характерная для многих современных наборов данных, приводит к явлению при котором традиционные алгоритмы, включая \kmeans{}, демонстрируют нестабильность и низкую точность результатов.
        Для преодоления этих ограничений активно применяются различные стратегии предварительной обработки данных, ключевую роль среди которых играет снижение размерности. Данный подход позволяет не только сократить вычислительную сложность алгоритмов, но и повысить устойчивость и точность кластеризации за счет устранения шумовых и коррелированных признаков.
        Несмотря на разнообразие существующих методов снижения размерности, на практике остается ряд открытых проблем. Во-первых, не существует универсального критерия выбора метода снижения размерности для конкретной задачи кластеризации. Во-вторых, влияние предварительного снижения размерности на итоговое качество кластеризации далеко не всегда очевидно и требует тщательного эмпирического анализа.
